\chapter{Verifica e convalida} % Main chapter title

\label{Chapter3} % For referencing the chapter elsewhere, use \ref{Chapter2} 

\section{Progettazione dei casi di test}

Questa sezione descrive nel dettaglio la progettazione dei casi di test e le motivazioni delle strategie adottate. L’obiettivo è definire un insieme di casi di test efficace, in grado di verificare che il sistema soddisfi i requisiti.

Tra i diversi approcci alla progettazione dei test, la strategia adottata si basa sul partizionamento in classi di equivalenza. Il dominio dei dati in ingresso viene suddiviso in classi, ciascuna composta da valori per cui il sistema dovrebbe comportarsi in modo equivalente, e per ogni classe si sceglie almeno un caso di test rappresentativo. Questa scelta riduce il numero di test necessari senza rinunciare alla copertura dei comportamenti distinti del sistema.

\section{Descrizione dei casi di testi in JUnit}
La presente sezione si occuperà di rappresentare i casi di test definiti per esaminare il comportamento delle funzionalità più complesse e per accertarne il corretto funzionamento. A questo scopo, è stato utilizzato il framework JUnit 5. \\
\indent Per facilitare la leggibilità, i casi di test sono stati suddivisi in base al sottosistema di appartenenza. I primi casi di test che verranno illustrati sono quelli definiti nella parte front-end; successivamente verranno descritti quelli della parte back-end.

I metodi selezionati (lato front-end) per gli Unit Test sono i seguenti:

\begin{quote}
\begin{enumerate}[label=\textit{\arabic*}]
   \item \texttt{selectToggleFilter()}: imposta un filtro di ricerca a partire 
   dai campi impostati nell’interfaccia grafica. Lo scopo è verificare che ogni combinazione di scelte, compresa l’assenza di selezione, produca un filtro valido.
  
  \item \texttt{setDate()}: traduce una data in formato UTC in una data leggibile dall'utente. 
  Lo scopo è verificare che le date significative (passaggio di mezzanotte o di fine anno, valore mancante o malformato) 
  siano convertite correttamente.
\end{enumerate} 
\end{quote}

\noindent Per la definizione delle classi di equivalenza è stato adottato l'\textit{approccio basato sulla funzionalità}, 
in cui il processo di partizionamento si basa rigorosamente sui requisiti
e non sull'implementazione del codice. \\
\indent Le classi di equivalenza sono pensate in relazione al
comportamento del metodo. Quindi i parametri non sono stati valutati in modo isolato, ma si
è tenuto conto del contesto e delle relazioni logiche esistenti tra di essi.

\subsection{Metodo \texttt{selectToggleFilter()}}

\noindent La firma del metodo è la seguente:

\begin{minipage}{\linewidth}
\centering\begin{minted}{java}
void selectToggleFilter(ToggleGroup group, String filterValue)
\end{minted}
\end{minipage}

\noindent\texttt{filterValue} rappresenta la tipologia di filtro selezionata appartenente ad uno specifico
gruppo, il quale è stato indicato nel primo parametro, \texttt{group}.

Per il parametro \texttt{group} il dominio è 
$G = \left\{ \text{typologyGroup}, \text{stateGroup},\text{priorityGroup}\right\}$ 
I gruppi vengono creati nativamente dall'interfaccia tramite FXML, venendo 
già popolati con i rispettivi \texttt{RadioButton}. Inoltre, il riferimento
non è mai nullo e si ottiene la seguente partizione

\noindent$$C_{1} =\left\{\text{typologyGroup}\right\},
C_{2} =\left\{\text{stateGroup}\right\},
C_{3} =\left\{\text{priorityGroup}\right\}$$

\noindent Per il parametro \texttt{filterValue}, sia $S = \left\{ w|w \text{ rappresenta una stringa} \right\}\cup \left\{ \textit{NIL} \right\}$  
l'insieme di tutte le stringhe incluso il valore nullo, sono state identificate
le seguenti classi di equivalenza

\noindent$$C_{4} =\left\{\text{BUG},\text{FEATURE},\text{DOCUMENTATION},\text{QUESTION}\right\}$$
$$C_{5} =\left\{\text{TODO},\text{CHIUSO}\right\},C_{6} =\left\{\text{ALTA},\text{MEDIA},\text{BASSA}\right\}$$
$$C_{7} = S\smallsetminus \bigcup_{i=4}^{6}C_{i}$$

\noindent Secondo l'implementazione del metodo \texttt{selectToggleFilter}, ogni classe è costituita da valori validi.
Per quanto riguarda la classe $C_7$, se la stringa non corrisponde a nessun \texttt{UserData} dei 
\texttt{RadioButton} definiti, il metodo lascia il gruppo senza selezione.

\noindent Il criterio di copertura adottato \textit{Each Choice Coverage}, che ha permesso
di generare la seguente batteria di test

\begin{table}[ht]
  \centering
  \small
  \caption{Casi di test}
  \label{tab:test-selecttogglefilter}
  \begin{tabular}{@{}llllp{4cm}@{}}
    \toprule
    Caso & \texttt{group} & \texttt{filterValue} & Classi & Esito atteso \\
    \midrule
    TC1 & \texttt{typologyGroup} & \texttt{BUG}  & $C_1, C_4$ & pulsante BUG selezionato  \\
    TC2 & \texttt{priorityGroup} & \texttt{ALTA} & $C_3, C_6$ & pulsante ALTA selezionato \\
    TC3 & \texttt{stateGroup}    & \texttt{TODO} & $C_2, C_5$ & pulsante TODO selezionato \\
    TC4 & \texttt{stateGroup} & \texttt{BLOCCATO}  & $C_2, C_7$ & nessuna selezione         \\
    \bottomrule
  \end{tabular}
\end{table}

\clearpage

\begin{listing}[!ht]
\begin{minted}{java}
@Test
    void testFilterTypologyValid()  {
        FilterFake view = new FilterFake();
        ToggleGroup typologyGroup = new ToggleGroup();

        view.toggle("BUG",
                typologyGroup);
        view.toggle("FEATURE",
                typologyGroup);
        view.toggle("DOCUMENTATION",
                typologyGroup);
        view.toggle("QUESTION",
                typologyGroup);

        view.selectToggleGroup(typologyGroup,
                "BUG");

        assertNotNull(typologyGroup
                .getSelectedToggle());
        assertEquals("BUG", typologyGroup
                .getSelectedToggle()
                .getUserData());
    }

    @Test
    void testFilterPriorityValid()  {
        FilterFake view = new FilterFake();
        ToggleGroup priorityGroup = new ToggleGroup();

        view.toggle("ALTA",
                priorityGroup);
        view.toggle("MEDIA",
                priorityGroup);
        view.toggle("BASSA",
                priorityGroup);

        view.selectToggleGroup(priorityGroup,
                "ALTA");

        assertNotNull(priorityGroup
                .getSelectedToggle());
        assertEquals("ALTA", priorityGroup
                .getSelectedToggle()
                .getUserData());
    }
\end{minted}
\end{listing}

\clearpage

\begin{listing}[!ht]
\begin{minted}{java}
    @Test
    void testFilterStateValid()  {
        FilterFake view = new FilterFake();
        ToggleGroup stateGroup = new ToggleGroup();

        view.toggle("TODO",
                stateGroup);
        view.toggle("CHIUSO",
                stateGroup);

        view.selectToggleGroup(stateGroup,
                "TODO");

        assertNotNull(stateGroup
                .getSelectedToggle());
        assertEquals("TODO", stateGroup
                .getSelectedToggle()
                .getUserData());
    }

    @Test
    void testFilterBadValid()  {
        FilterFake view = new FilterFake();
        ToggleGroup stateGroup = new ToggleGroup();

        view.toggle("TODO",
                stateGroup);
        view.toggle("CHIUSO",
                stateGroup);

        view.selectToggleGroup(stateGroup,
                "APERTO");
        assertNull(stateGroup
                .getSelectedToggle());
    }
\end{minted}
\end{listing}

\begin{listing}[!ht]
\begin{minted}{java}

\end{minted}
\end{listing}

\subsection{Metodo \texttt{setDate()}}
Poiché non sono presenti metodi pubblici non banali con due o 
più parametri che non effettuino chiamate verso il back-end, si 
è scelto di testare uno dei metodi ritenuti più complessi, setDate, 
pur accettando un solo parametro.

\noindent La firma del metodo è la seguente:

\begin{minipage}{\linewidth}
\centering\begin{minted}{java}
void setDate(LocalDateTime date)
\end{minted}
\end{minipage}

Il metodo presenta un dominio di input molto ampio. In base al 
comportamento atteso, in cui ad esempio intervalli inferiori a un minuto 
vengono tutti indicati come “\textit{Just now}”, è stata ricavata la seguente partizione:

\noindent$$C_{1} =\left\{NIL\right\},
C_2 = \left\{ w \mid w \text{ rappresenta una data del futuro}\right\},$$
$$C_3 = \left\{ \Delta \mid 0 \le \Delta < 60\text{s} \right\}$$
dove $\Delta$ rappresenta il tempo trascorso tra la data fornita e il tempo attuale.
$$C_4 = \left\{ \Delta \mid 0 < \Delta < 59\text{min} \right\},
C_5 = \left\{ \Delta \mid 0 <  \Delta < 24\text{h} \right\},
C_6 = \left\{ \Delta \mid 0 <  \Delta < 7\text{d} \right\}$$
$$C_7 = \left\{ \Delta \mid 7d \le \Delta < 1\text{m} \right\},
C_8 = \left\{ \Delta \mid 0 < \Delta < 12\text{m} \right\},$$
$$C_9 = \left\{ \Delta \mid \Delta > 0 \right\}$$

dove 1\text{m} rappresenta 1 mese, 7d rappresenta 7 giorni e l'ultima classe di equivalenza
esprime l'intervallo in anni. \\
\indent Il criterio di copertura adottato Each Choice Coverage, che ha permesso di generare la
seguente batteria di test.


\begin{listing}[!ht]
\begin{minted}{java}
    @Test
    void testDateNull()  {
        assertThrows(NullPointerException.class,
                () -> {
            NoticeDate notice =
                    new NoticeDate();
            notice.setDate(null);
        });
    }

    @Test
    void testPastDateYear()  {
        NoticeDate view = new NoticeDate();
        view.setDate(LocalDateTime
                .now(ZoneOffset.UTC)
                .minusYears(3));

        assertNotNull(view);
        assertEquals("3 years ago", view
                .getDate());
    }

    @Test
    void testPastDateMonth()  {
        NoticeDate view = new NoticeDate();
        view.setDate(LocalDateTime
                .now(ZoneOffset.UTC)
                .minusMonths(11));

        assertNotNull(view);
        assertEquals("11 months ago", view
                .getDate());
    }

\end{minted}
\end{listing}

\begin{listing}[!ht]
\begin{minted}{java}
    @Test
    void testPastDateWeek()  {
        NoticeDate view = new NoticeDate();
        view.setDate(LocalDateTime
                .now(ZoneOffset.UTC)
                .minusDays(14));

        assertNotNull(view);
        assertEquals("2 weeks ago", view
                .getDate());
    }

    @Test
    void testPastDateDay()  {
        NoticeDate view = new NoticeDate();
        view.setDate(LocalDateTime
                .now(ZoneOffset.UTC)
                .minusDays(6));

        assertNotNull(view);
        assertEquals("6 days ago", view
                .getDate());
    }

    @Test
    void testPastDateHour()  {
        NoticeDate view = new NoticeDate();
        view.setDate(LocalDateTime
                .now(ZoneOffset.UTC)
                .minusHours(23));

        assertNotNull(view);
        assertEquals("23 hours ago", view
                .getDate());
    }

    @Test
    void testPastDateMinute()  {
        NoticeDate view = new NoticeDate();
        view.setDate(LocalDateTime
                .now(ZoneOffset.UTC)
                .minusMinutes(59));

        assertNotNull(view);
        assertEquals("59 minutes ago", view
                .getDate());
    }

    @Test
    void testPastDateSecond()  {
        NoticeDate view = new NoticeDate();
        view.setDate(LocalDateTime
                .now(ZoneOffset.UTC)
                .minusSeconds(5));

        assertNotNull(view);
        assertEquals("Just now", view
                .getDate());
    }

\end{minted}
\end{listing}

\clearpage

\begin{listing}[!ht]
\begin{minted}{java}
    @Test
    void testFutureDate()  {
        NoticeDate view = new NoticeDate();
        view.setDate(LocalDateTime
                .now(ZoneOffset.UTC)
                .plusMonths(5));

        assertNotNull(view);
        assertEquals("Just now", view
                .getDate());
    }
\end{minted}
\end{listing}

\begin{table}[ht]
  \centering
  \small
  \caption{Casi di test}
  \label{tab:test-selecttogglefilter}
  \begin{tabular}{@{}llllp{4cm}@{}}
    \toprule
    Caso & \texttt{date} & Classi & Esito atteso \\
    \midrule
    TC1 & \texttt{null} & $C_1$ & \texttt{NullPointerException}  \\
    TC2 & \texttt{now() - 3 anni} & $C_9$ & "3 years ago" \\
    TC3 & \texttt{now() - 11 mesi}    & $C_8$ & "11 months ago" \\
    TC4 & \texttt{now() - 2 settimane}    & $C_7$ & "2 weeks ago"       \\
    TC5 & \texttt{now() - 6 giorni} & $C_6$ & "6 days ago"  \\
    TC6 & \texttt{now() - 23 ore} & $C_5$ & "23 hours ago" \\
    TC7 & \texttt{now() - 59 minuti}    & $C_4$ & "59 minutes ago" \\
    TC8 & \texttt{now() - 5 secondi}    & $C_3$ & "Just now"       \\
    TC9 & \texttt{now() + 5 secondi}    & $C_2$ & "Just now"       \\

    \bottomrule
  \end{tabular}
\end{table}

\section{Testing lato back-end}
\subsection{Test plan: creazione di una issue}

\subsubsection{Obiettivo} 

Verificare che un utente interno (amministratore o utente normale) possa 
segnalare un’issue e che il sistema applichi le seguenti regole di servizio: 
titolo e descrizione obbligatori, stato iniziale TODO, priorità NONE se non 
indicata e creazione di un oggetto bug solo per le issue di tipo BUG.  
Verificare inoltre che gli utenti non autorizzati vengano respinti. 

\subsubsection{Attori del test} 

- L'endpoint POST /api/issues (che corrisponde a IssueController.creaIssue), 
che controlla il ruolo dell'utente e la validità dei dati ricevuti. 
- La classe CreaIssueRequest, che contiene i vincoli sui dati in ingresso 
(campi obbligatori e lunghezze). 
- Il metodo UtenteInternoService.creaIssue, che applica le regole di creazione. 

\subsubsection{Ambiente di testing} 

Back-end e database MySQL avviati ed operativi; nel database sono 
presenti l'amministratore di default, un utente normale e uno stakeholder.  

Prima di ogni caso di test l'utente indicato ha effettuato il login
 e la richiesta contiene il suo cookie di sessione.  

Tranne dove indicato, il corpo della richiesta è sempre: 

{ “titolo”: “Login non funziona”, “descrizione”: “Errore 500 al login”, “tipo”: “BUG”, 
“priorita”: “ALTA”, “immagine”: null } 

\clearpage

\begin{table}[htbp]
  \centering
  \renewcommand{\arraystretch}{1.4}
  \begin{tabular}{|>{\raggedright\arraybackslash}p{0.6cm}|>{\raggedright\arraybackslash}p{3.6cm}|>{\raggedright\arraybackslash}p{2.6cm}|>{\raggedright\arraybackslash}p{3.2cm}|>{\raggedright\arraybackslash}p{2.6cm}|}
    \hline
    \textit{Id} & \textit{Precondizioni} & \textit{Input} & \textit{Output effettivo} & \textit{Postcondizioni} \\
    \hline
    TP1 & Nessun login & Corpo di riferimento & 401 Unauthorized, ``Utente non autenticato'' & Nessuna issue creata \\
    \hline
    TP2 & Login come stakeholder & Corpo di riferimento & 403 Forbidden, ruolo non ammesso, ``Operazione non permessa'' & Nessuna issue creata \\
    \hline
    TP3 & Login come utente normale & Corpo di riferimento & 201 Created & Bug creato (nel db righe in Issues e in Bugs). UtenteCreatore = l'utente loggato\\
    \hline
    TP4 & Login come amministratore & Corpo di riferimento & 201& Bug creato (nel db righe in Issues e in Bugs). UtenteCreatore = l'utente loggato \\
    \hline
    TP5 & Login come utente normale & tipo = FEATURE & 201& Riga creata solo in Issues \\
    \hline
    TP6 & Login come utente normale & priorità assente & 201, priorità NONE& Issue creata \\
    \hline
    TP7 & Login come utente normale & titolo assente & 400 Bad request, richiesta non valida, “Il titolo è obbligatorio”
& Nessuna issue creata\\
    \hline
    TP8 & Login come utente normale & titolo di 1 carattere & 201 & Issue creata \\
    \hline
    TP9 & Login come utente normale & titolo di 100 caratteri & 201 & Issue creata\\
    \hline
    TP10 & Login come utente normale & titolo di 101 caratteri & 400 , “Il titolo può avere al massimo 100 caratteri” & Nessuna issue creata \\
    \hline
    TP11 & Login come utente normale& descrizione di 500 caratteri & 201 & Issue creata \\
    \hline
    TP12 & Login come utente normale & descrizione di 501 caratteri & 400, “La descrizione può avere al massimo 500 caratteri” & Nessuna issue creata \\
    \hline
    TP13 & Login come utente normale& tipo assente & 400, “Il tipo è obbligatorio” & Nessuna issue creata \\
    \hline
    TP14 & Login come utente normale& tipo = "PIPPO" oppure priorità = "URGENTE"& 400 & Nessuna issue creata\\
    \hline
    TP15 & Login come utente normale& immagine di 2 MB& 201, con immagine & Issue creata\\
    \hline
    TP16 & Login come utente normale& JSON non valido& 400, “Richiesta non valida” & Nessuna issue creata\\
    \hline
  \end{tabular}
\end{table}

\clearpage

\subsubsection{Criteri di superamento}
Il test plan è superato se tutti i casi di test producono l'output e 
le postcondizioni attese.

\subsection{Metodo \texttt{creaIssue()} }
La firma del metodo è la seguente: 

\begin{minipage}{\linewidth}
\centering\begin{minted}{java}
creaIssue(idCreatore, titolo, descrizione, tipo, priorita, immagine) 
\end{minted}
\end{minipage}

Da questo metodo abbiamo identificato le seguenti classi di equivalenza: 

\noindent idCreatore: CE1 = $\left\{ x \text{ è l'id di un utente interno esistente} \right\}$ (Valida) \\
idCreatore: CE2 = $\left\{ x \text{ non corrisponde a nessun utente} \right\}$  (Non valida) \\
tipo: CE3 = $\left\{ x = \text{ BUG} \right\}$ (Valida) \\
tipo: CE4 = $\left\{ x \in \left\{ QUESTION, DOCUMENTATION, FEATURE \right\} \right\}$ (Valida) \\
priorita: CE5 =  $\left\{ x \in \left\{ BASSA, MEDIA, ALTA \right\} \right\}$ (Valida) \\
priorita: CE6 = $\left\{ x \text{ è NULL} \right\}$ (Valida, la priorità diventa NONE) \\
immagine: CE7 =  $\left\{ x \text{ è un array di byte} \right\}$ (Valida) \\
immagine: CE8 = $\left\{ x \text{ è NULL} \right\}$ (Valida, issue senza immagine) \\
titolo: CE9 = $\left\{ x \text{ è una stringa non vuota di al massimo 100 caratteri} \right\}$ (Valida) \\
descrizione: CE10 = $\left\{ x \text{ è una stringa non vuota di al massimo 500 caratteri} \right\}$ (Valida) \\

\noindent Le classi non valide di titolo e descrizione ($\left\{ x \text{ è NULL oppure ""} \right\}$ e \\
$\left\{ x \text{ supera la lunghezza massima} \right\}$) non 
sono state considerate: la validazione della richiesta le rifiuta con errore 400 prima che creaIssue venga chiamato.  


\begin{listing}[!ht]
\begin{minted}{java}

class UtenteInternoServiceTest {

    private final Integer idUtenteCreatore = 7;

    private IssueRepository issueRepository;
    private UtenteInternoRepository utenteInternoRepository;
    private UtenteInternoService utenteInternoService;
    private UtenteNormale utenteCreatore;

    // Setup: crea i mock per le repository e la classe contenente 
    // creaIssue, infine crea un utente con id=7
    @BeforeEach
    void setUp() {
        issueRepository = mock(IssueRepository.class);
        utenteInternoRepository = mock(UtenteInternoRepository.class);
        utenteInternoService = new UtenteInternoService(issueRepository, 
        utenteInternoRepository);

        utenteCreatore = new UtenteNormale();
        utenteCreatore.setId(idUtenteCreatore);
        when(utenteInternoRepository.findById(idUtenteCreatore))
        .thenReturn(Optional.of(utenteCreatore));
    }

\end{minted}
\end{listing}

\clearpage

\begin{listing}[!ht]
\begin{minted}{java}

 // CI1 - Utente creatore inesistente: la repository non 
     // lo trova e lancia l'eccezione
    @Test
    void creaIssue_creatoreNonTrovato_lanciaEccezione() {
        // precondizione: l'id 99 non corrisponde a nessun utente
        when(utenteInternoRepository.findById(99))
        .thenReturn(Optional.empty());

        // si lancia l'eccezione con questo id
        assertThrows(EntityNotFoundException.class, () ->
                        utenteInternoService
                        .creaIssue(99, "Titolo", "Descrizione",
                                TipoIssue.BUG, Priorita.ALTA, null));

        // controlla che nessuna issue è stata salvata
        verify(issueRepository, times(0)).save(any());
    }

     // CI2 - Issue tipo BUG: viene creato un Bug, non 
     // assegnato e non duplicato.
    @Test
    void creaIssue_tipoBug_creaBug() {

        utenteInternoService.creaIssue(idUtenteCreatore, "Login non funziona", 
        "Errore 500 al login",
                TipoIssue.BUG, Priorita.ALTA, null);

        // assert
        verify(issueRepository).save(any(Bug.class));
        verify(issueRepository).save(argThat(allOf(
                hasProperty("tipo", equalTo(TipoIssue.BUG)),
                hasProperty("assegnatario", is(nullValue())),
                hasProperty("duplicato", is(false)))));
    }

     // CI3 — Issue tipo QUESTION
    @Test
    void creaIssue_tipoQuestion() {

        utenteInternoService.creaIssue(idUtenteCreatore, "Avete provato 
        la nuova funzione dell'IDE?",
                "Vorrei qualche delucidazione in merito",
                TipoIssue.QUESTION, Priorita.BASSA, null);

        // assert: salvata un'issue di tipo QUESTION
        verify(issueRepository).save(argThat(hasProperty("tipo", 
        equalTo(TipoIssue.QUESTION))));
    }

\end{minted}
\end{listing}

\clearpage

\begin{listing}[!ht]
\begin{minted}{java}

// CI4 - Issue tipo DOCUMENTATION
    @Test
    void creaIssue_tipoDocumentation() {

        utenteInternoService.creaIssue(idUtenteCreatore, "Cosa restituisce 
                questo endpoint?",
                "Manca una documentazione adeguata",
                TipoIssue.DOCUMENTATION, Priorita.MEDIA, null);

        // assert
        verify(issueRepository).save(argThat(hasProperty("tipo", 
        sequalTo(TipoIssue.DOCUMENTATION))));
    }

 // CI5 - Issue tipo FEATURE
    @Test
    void creaIssue_tipoFeature() {

        utenteInternoService.creaIssue(idUtenteCreatore, 
                "Dimenticato di implementare la crittografia dei dati",
                "Feature importante richiesta dal cliente",
                TipoIssue.FEATURE, Priorita.ALTA, null);

        // assert
        verify(issueRepository).save(argThat(hasProperty("tipo", 
        equalTo(TipoIssue.FEATURE))));
    }

    // CI6 - Issue con priorità non indicata: vale NONE.
    @Test
    void creaIssue_prioritaNull() {

        utenteInternoService.creaIssue(idUtenteCreatore, 
        "Login non funziona", "Errore 500 al login",
                TipoIssue.FEATURE, null, null);

        // assert
        verify(issueRepository)
        .save(argThat(hasProperty("priorita", 
        equalTo(Priorita.NONE))));
    }

    // CI7 - Issue con priorità indicata
    @Test
    void creaIssue_prioritaIndicata() {
        utenteInternoService.creaIssue(idUtenteCreatore, 
        "Login non funziona", 
        "Errore 500 al login",
                TipoIssue.FEATURE, Priorita.MEDIA, null);

        // assert
        verify(issueRepository)
        .save(argThat(hasProperty("priorita", 
        equalTo(Priorita.MEDIA))));
    }
   
\end{minted}
\end{listing}

\clearpage

\begin{listing}[!ht]
\begin{minted}{java}
    // CI8 - Regole di servizio valide per ogni issue: imposta data 
    // di creazione e stato=TODO
    @Test
    void creaIssue_impostaStatoTodoEData() {

        utenteInternoService.creaIssue(idUtenteCreatore, 
        "Login non funziona", "Errore 500 al login",
                TipoIssue.QUESTION, Priorita.ALTA, null);

        // assert
        verify(issueRepository).save(argThat(allOf(
                hasProperty("stato", equalTo(StatoIssue.TODO)),
                hasProperty("data", is(not(nullValue()))),
                hasProperty("utenteCreatore", equalTo(utenteCreatore)),
                hasProperty("titolo", equalTo("Login non funziona")),
                hasProperty("descrizione", equalTo("Errore 500 al login")),
                hasProperty("immagine", is(nullValue())))));
    }

  // CI9 - Issue con immagine presente
    @Test
    void creaIssue_conImmagine() {
        byte[] immagine = {1, 2, 3, 4};

        utenteInternoService.creaIssue(idUtenteCreatore, 
        "Login non funziona", "Errore 500 al login",
                TipoIssue.BUG, Priorita.ALTA, immagine);

        // assert
        verify(issueRepository).save(argThat(hasProperty("immagine",
         equalTo(immagine))));
    }
}
\end{minted}
\end{listing}

\noindent CI1 (riga 57): creaIssue(99, "Titolo", "Descrizione", BUG, ALTA, NIL), copre CE2, CE3, CE5, CE8, CE9, CE10 \\
CI2 (riga 72): creaIssue(7, "Login non funziona", "Errore 500 al login", BUG, ALTA, NIL), copre CE1, CE3, CE5, CE8, CE9, CE10 \\
CI3 (riga 87): creaIssue(7, "Avete provato la nuova funzione dell'IDE?", "Vorrei qualche delucidazione in merito", QUESTION, BASSA, NIL), copre CE1, CE4, CE5, CE8, CE9, CE10 \\
CI4 (riga 99): creaIssue(7, "Cosa restituisce questo endpoint?", "Manca una documentazione adeguata", DOCUMENTATION, MEDIA, NIL), copre CE1, CE4, CE5, CE8, CE9, CE10 \\
CI5 (riga 111): creaIssue(7, "Dimenticato di implementare la crittografia dei dati", "Feature importante richiesta dal cliente", FEATURE, ALTA, NIL), copre CE1, CE4, CE5, CE8, CE9, CE10 \\
CI6 (riga 123): creaIssue(7, "Login non funziona", "Errore 500 al login", FEATURE, NIL, NIL), copre CE1, CE4, CE6, CE8, CE9, CE10 \\
CI7 (riga 134): creaIssue(7, "Login non funziona", "Errore 500 al login", FEATURE, MEDIA, NIL), copre CE1, CE4, CE5, CE8, CE9, CE10 \\
CI8 (riga 144): creaIssue(7, "Login non funziona", "Errore 500 al login", QUESTION, ALTA, NIL), copre CE1, CE4, CE5, CE8, CE9, CE10 \\
CI9 (riga 161): creaIssue(7, "Login non funziona", "Errore 500 al login", BUG, ALTA, {1, 2, 3, 4}), copre CE1, CE3, CE5, CE7, CE9, CE10 \\

\subsubsection{Criterio di copertura delle classi di equivalenza}
Come criterio di copertura si è adottato l'Each Choice Coverage (ECC): ogni 
classe di equivalenza di ogni parametro è rappresentata da almeno un caso di 
test. I test da CI1 a CI9 coprono così tutte le classi individuate.
